\section*{الفصل \ref{ch:b3:statistical-thermo-applied} --- تطبيقات الترموديناميك الإحصائي}

\begin{solution}{exo:b3:statistical-thermo-applied:1}
$\frac32R + R + (3 \times 3 - 5)R = 6.5R = \qty{54.0}{J.K^{-1}.mol^{-1}}$. في درجة حرارة الغرفة يكون الانسحاب
والدوران كلاسيكيين، والتمددان مجمّدين، والانحناء المنحل
مرتين، وهو أدنى الاهتزازات، مثارًا جزئيًا.
\end{solution}

\begin{solution}{exo:b3:statistical-thermo-applied:2}
$x = 1$: $\eu/(\eu - 1)^2 = 0.921$. $x = 10$: $100\eu^{10}/(\eu^{10} - 1)^2 \approx 100\eu^{-10} = \num{4.5e-3}$.
\end{solution}

\begin{solution}{exo:b3:statistical-thermo-applied:3}
\qty{20}{K}: $x_{\mathrm{para}} = 1/(1 + 3 \times 3\eu^{-2 \times 85.35/20}) = 1/(1 + 9 \times \num{1.96e-4}) = 0.998$.
\qty{77}{K}: المجموع الزوجي $1 + 5\eu^{-6.651} = 1.0065$؛ والمجموع الفردي $3(3\eu^{-2.217} + 7\eu^{-13.30}) = 0.9806$؛
$x_{\mathrm{para}} = 1.0065/1.9871 = 0.51$.
\end{solution}

\begin{solution}{exo:b3:statistical-thermo-applied:4}
من أجل $I = 1$ توجد $3 \times 2 = 6$ حالات سبين متناظرة و$3 \times 1 = 3$ ضد متناظرة؛
والديوترونات بوزونات، فالدالة الموجية الكلية متناظرة: حالات السبين المتناظرة
مع $J$ الزوجية (أورثو، الوزن 6)، وضد المتناظرة مع $J$ الفردية (بارا، الوزن
3). والدوتيريوم الأورثو، الذي يحتوي $J = 0$، مستقر عند 0 K؛ وفي درجة حرارة الغرفة
أورثو:بارا = 2:1.
\end{solution}

\begin{solution}{exo:b3:statistical-thermo-applied:5}
$\Delta_rE_0 = 2 \times 1883.7 - 2170.3 - 1542.3 = \qty{54.8}{cm^{-1}}$؛ $\eu^{-54.8/207.2} = 0.768$، 
و$4 \times 0.768 = 3.07$. والقيمة الكاملة، 3.26، أعلى بنسبة 6\,\%: فالعوامل الانسحابية والدورانية
والاهتزازية لا تتلاشى بالضبط إلا في النهاية الكلاسيكية، 
ودوران هذه الجزيئات الخفيفة ما زال مكمّى جزئيًا عند \qty{298}{K}.
\end{solution}

\begin{solution}{exo:b3:statistical-thermo-applied:6}
\ce{N2O}: اتجاهان، $R\ln2 = \qty{5.76}{J.K^{-1}.mol^{-1}}$. \ce{CH3D}: أربعة، $R\ln4 =
\qty{11.53}{J.K^{-1}.mol^{-1}}$ (حدود عليا، تُبلغ إذا كانت الترتيبات عشوائية).
\end{solution}

\begin{solution}{exo:b3:statistical-thermo-applied:7}
$x = 797.6/300 = 2.659$: دالة أينشتاين $x^2\eu^{-x}/(1 - \eu^{-x})^2 = 0.572$؛ $C_V/R = 3.072$،
$C_p = 4.072R = \qty{33.86}{J.K^{-1}.mol^{-1}}$، أي أدنى بنسبة 0.4\,\% من قيمة الجدول (اللاتوافق).
\end{solution}

\begin{solution}{exo:b3:statistical-thermo-applied:8}
$x = 1.029$: $x^2\eu^{-x}/(1 - \eu^{-x})^2 = 0.916$، أي \qty{7.62}{J.K^{-1}.mol^{-1}}، وهي 92\,\% من $R$: فاهتزاز
اليود شبه كلاسيكي في درجة حرارة الغرفة.
\end{solution}

\begin{solution}{exo:b3:statistical-thermo-applied:9}
B أغنى بالنظير \ce{^{18}O}. $\alpha_{\mathrm{A/B}} = (1 - 0.0100)/(1 + 0.0020) = 0.98802$.
\end{solution}

\begin{solution}{exo:b3:statistical-thermo-applied:10}
$\Delta\mathrm{ZPE} = \frac12(1 - 1/\sqrt2)(2000 - 3000) = \qty{-146}{cm^{-1}}$؛ $K = \eu^{146.4/207.2} = 2.0$.
يذهب الدوتيريوم تفضيلًا إلى X، الرابطة الأصلب، حيث يكون اقتصاده في \omterm{def:b3:quantum-model-systems:oscillator}{طاقة
النقطة الصفرية} أكبر.
\end{solution}

\begin{solution}{exo:b3:statistical-thermo-applied:11}
$\Delta_rH^\circ = R\ln10\,(-1.766 + 3.392)/(1/900 - 1/1100) = 19.145 \times 1.626/\num{2.020e-4} =
\qty{154}{kJ/mol}$. تحمل الذرتان $2 \times \frac52RT$ من الأنتالبي؛ ويحمل الجزيء $\frac52RT + RT$
(الدوران) مضافًا إليها طاقته الاهتزازية $R\theta_{\mathrm V}/(\eu^{\theta_{\mathrm V}/T} - 1)$، الأقل قليلًا من $RT$:
والفرق، نحو \qty{5}{kJ/mol} عند \qty{1000}{K}، يضاف إلى $D_0$.
\end{solution}

\begin{solution}{exo:b3:statistical-thermo-applied:12}
المستويان 0 ($g = 1$) و$6hcB$ ($g = 5$): $q = 1 + 5\eu^{-x}$، $\langle\varepsilon\rangle/kT = 5x\eu^{-x}/q$، $\langle\varepsilon^2\rangle/(kT)^2 =
5x^2\eu^{-x}/q$، و$C/R = \langle\varepsilon^2\rangle/(kT)^2 - (\langle\varepsilon\rangle/kT)^2 = 5x^2\eu^{-x}/(1 + 5\eu^{-x})^2$. وعند
\qty{100}{K}، $x = 5.121$: $C/R = 0.783/1.061 = 0.738$.
\end{solution}

\begin{solution}{pb:b3:statistical-thermo-applied:1}
\textbf{1.} $2 \times 2 = 4$: ثلاث متناظرة (الثلاثية)، وواحدة ضد متناظرة (الأحادية).
\textbf{2.} البروتونات فرميونات، فالدالة الموجية الكلية تغيّر إشارتها
بالتبادل؛ والدالة الدورانية تعطي $(-1)^J$: فقيم $J$ الزوجية تقترن \omterm{def:b3:many-electron-atoms:singlet-triplet}{بالحالة
الأحادية} ضد المتناظرة (بارا)، وقيم $J$ الفردية \omterm{def:b3:many-electron-atoms:singlet-triplet}{بالحالة الثلاثية} المتناظرة (أورثو).
\textbf{3.} المستويات الزوجية 1، والفردية 3.
\textbf{4.} عند \qty{300}{K} ($T = 3.5\,\theta_{\mathrm R}$) يكاد المجموعان الدورانيان الزوجي والفردي يتساويان؛
وبترجيحهما بالوزنين 1 و3 يعطيان 1:3 (نسبة البارا 0.2507).
\textbf{5.} $F(1) = 2 \times 59.322 - 4 \times 0.0471 = \qty{118.46}{cm^{-1}} = \qty{1.417}{kJ/mol}$.
\textbf{6.} $1/(1 + 9\eu^{-170.4/20.37}) = 1/(1 + 9 \times \num{2.3e-4}) = 0.998$.
\textbf{7.} مع $J = 0$ و1 وحدهما، يعني النصف بارا $9\eu^{-2\theta_{\mathrm R}/T} = 1$، $T = 2\theta_{\mathrm R}/\ln9 =
0.91\,\theta_{\mathrm R} = \qty{78}{K}$: فالتنافس بين الوزن 9 للمستوى $J = 1$ 
وعامل بولتزمان له.
\textbf{8.} 0.75. فدون حفاز يحتاج التحول إلى أيام، والتسييل إلى
ساعات.
\textbf{9.} $(0.75 - 0.002) \times 1.417 = \qty{1.060}{kJ/mol}$.
\textbf{10.} $1.060/\num{2.016e-3} = \qty{526}{kJ/kg}$.
\textbf{11.} $0.905/\num{2.016e-3} = \qty{449}{kJ/kg}$.
\textbf{12.} حرارة التحول 1.17 مرة أنتالبي التبخر.
\textbf{13.} إنها تتولد داخل السائل، لا تدخل عبر الجدران.
\textbf{14.} $x_{\mathrm o}(t) = x_{\mathrm{eq}} + (0.75 - x_{\mathrm{eq}})\eu^{-kt}$، مع $x_{\mathrm{eq}} = 0.002$.
\textbf{15.} $0.002 + 0.748\eu^{-0.274} = 0.571$.
\textbf{16.} $(0.75 - 0.571) \times 1.417 = \qty{0.254}{kJ/mol}$.
\textbf{17.} $0.254/0.905 = 0.28$: أي 28\,\% من الخزان في يوم واحد.
\textbf{18.} بعد \qty{168}{h}، $x_{\mathrm o} = 0.112$، والحرارة \qty{0.904}{kJ/mol}، أي ما يساوي أنتالبي
التبخر: ففي هذا النموذج يفرغ الخزان. ويبالغ التقدير في
الفقد لأن الجزيئات المتبخرة تحمل معها \omterm{def:b3:statistical-thermo-applied:spin-isomers}{الهيدروجين الأورثو} غير المتحول؛
والفقد الحقيقي أصغر لكنه يبقى كبيرًا.
\textbf{19.} $t_{1/2} = \ln2/0.0114 = \qty{61}{h}$.
\textbf{20.} في جهاز التسييل، على طبقات حفاز عند درجات حرارة متتالية، كي
يزيل جهاز التبريد حرارة التحول قبل تخزين السائل.
\textbf{21.} نعم: نسبة البارا عند التوازن عند \qty{300}{K} هي 0.2507.
\textbf{22.} في الهيدروجين العادي لا تستطيع الصورتان أن تتحولا إحداهما إلى الأخرى، فيتجمد
دوران كل منهما على سلّم مستوياته الخاص؛ أما في الهيدروجين في حالة التوازن فيحوّل جزء من
الحرارة المقدَّمة البارا إلى أورثو.
\textbf{23.} نحو 3.57 قرب \qty{50}{K}: فالحرارة، إلى جانب إثارة الدوران، تحوّل
جزيئات من $J = 0$ إلى $J = 1$، الأعلى بمقدار \qty{118}{cm^{-1}}، وهو تفاعل أنتالبيه
موجب.
\textbf{24.} 1.50: الدوران مجمّد في الصورتين (البارا في $J = 0$، والأورثو في $J = 1$).
\textbf{25.} \textbf{حرارة التحول / أنتالبي التبخر $\approx 1.2$ للهيدروجين
العادي}: فتحوله وحده قد يبخر الخزان كله، ولهذا
يُحوَّل الهيدروجين إلى بارا قبل التخزين.
\end{solution}
